%%%%%%%%%%%%%%%%%%%%%%%%%%%%%%%%%%%%
\chapter{Contribution}
\label{ch:contribution}
% ---------------------------------------------------------------------------
% Section map (thesis_writeup/01_thesis_structure_v2.md §3):
%   3.1 Overview of the study                       [outline]
%   3.2 Data                                        [v0 DRAFT written overnight 17->18 Aug — for review]
%   3.3 The foundation model and its adaptation     [outline]
%   3.4 The experimental programme, dated           [outline; Table 3.4]
%   3.5 Hand-crafted ECG features                   [outline]
%   3.6 Analysis pipelines                          [outline; Table 3.6]
%   3.7 The informativeness-fidelity audit          [outline]
%   3.8 Transport, prediction and repair protocols  [outline]
% Every setup number in 3.2 traces to notes/dossiers/A_data_and_setup.md (A§x) and, through it,
% to the manifest / JSON / script line quoted there.
% ---------------------------------------------------------------------------

\section{Overview of the study}
\label{sec:overview}
\todo{3.1 — one page: design in one paragraph; Fig 3.1 pipeline schematic; novelty fences A–F with the ``to our knowledge'' caveat}

\section{Data}
\label{sec:data}

Two ECG cohorts are used throughout: a synthetic cohort generated by a cardiac digital twin, \medalcare, and a real clinical cohort, \ptbxl. Both are public. This section fixes the counts, label spaces and pre-processing conventions that every later chapter relies on; the four confirmed data-handling defects found in the August audit and their fixes are described where they belong (\S\ref{sec:preproc}), and their consequences are quantified in Chapter~\ref{ch:results}.

\subsection{MedalCare-XL: the synthetic cohort}
\label{sec:data-medalcare}

\medalcare~\citep{gillette2023medalcare} is a set of 12-lead ECGs computed from anatomically and electrophysiologically parameterised whole-heart simulations. The release used here contains 16,839 ten-second recordings sampled at 500\,Hz in standard lead order (I, II, III, aVR, aVL, aVF, V1--V6), each carrying exactly one of eight pathology labels: sinus rhythm, myocardial infarction (MI), right and left bundle-branch block, left atrial enlargement, interatrial conduction block, fibrotic atrial cardiomyopathy, and AV block. Of the three signal variants shipped per simulation (raw, noise-augmented, filtered) only the filtered variant (0.5--150\,Hz Butterworth) is used. \tracenote{A§1.1.1; MedalCare-XL/WP2_largeDataset_Noise/README.txt; scripts/prepare_medalcare.py:112}

The dataset's own training/validation/test partition (${\approx}70/15/15$) is defined by the anatomical model with which the QRST complexes were simulated, so that ECGs from the same anatomy---but different electrophysiological parameters---fall into one partition only. This project inherits that partition unchanged: the working manifest reads the split off the release's directory names, and agreement with the directory structure is 100\,\% (16,839/16,839 rows). \tracenote{A§1.1.5, A§7.1; scripts/add_medalcare_splits.py; data/medalcare_filtered_manifest_dataset_split.csv} Table~\ref{tab:medalcare} gives the resulting counts. The MI subset---7,797 recordings across eight occlusion folders that combine the culprit artery (LAD, LCX, RCA), a transmurality parameter ($0.3$ or $1.0$) and, for LCX, an anterior or posterior sub-location---is the cohort on which every $\thetavec$-decoding and territory experiment is run. \tracenote{A§1.1.4; data/theta_mi_build_summary.json→territory_8c_counts}

\begin{table}[tb]
\centering\small
\caption{\medalcare\ cohort as used in this thesis (filtered signal variant; the release's own model-disjoint partition). The eight-class labels are mutually exclusive. The MI rows are further split by occlusion folder; the four-class territory used downstream is derived from the simulated ischaemia angle $\phiang$, not from the folder name (\S\ref{sec:theta}).}
\label{tab:medalcare}
\begin{tabular}{lrrrr}
\toprule
 & train & val & test & total \\
\midrule
sinus rhythm & 900 & 200 & 200 & 1,300 \\
myocardial infarction (MI) & 5,347 & 1,250 & 1,200 & 7,797 \\
right bundle-branch block & 898 & 200 & 200 & 1,298 \\
left bundle-branch block & 900 & 200 & 200 & 1,300 \\
left atrial enlargement & 1,040 & 130 & 130 & 1,300 \\
interatrial conduction block & 994 & 124 & 126 & 1,244 \\
fibrotic atrial cardiomyopathy & 1,040 & 130 & 130 & 1,300 \\
AV block & 900 & 200 & 200 & 1,300 \\
\midrule
all recordings & 12,019 & 2,434 & 2,386 & 16,839 \\
\midrule
\multicolumn{5}{l}{\emph{MI subset by occlusion folder (artery\_transmurality[\_sub-location])}}\\
LAD\_0.3 / LAD\_1.0 & 899 / 900 & 200 / 200 & 200 / 200 & 1,299 / 1,300 \\
LCX\_0.3\_ant / LCX\_0.3\_post & 450 / 450 & 150 / 100 & 100 / 100 & 700 / 650 \\
LCX\_1.0\_ant / LCX\_1.0\_post & 400 / 450 & 100 / 100 & 100 / 100 & 600 / 650 \\
RCA\_0.3 / RCA\_1.0 & 898 / 900 & 200 / 200 & 200 / 200 & 1,298 / 1,300 \\
\midrule
\multicolumn{5}{l}{\emph{MI subset by $\phiang$-derived four-class territory}}\\
Anteroseptal (\AS) & 1,799 & 400 & 400 & 2,599 \\
Anterolateral (\AL) & 1,300 & 350 & 300 & 1,950 \\
Inferior (\INF) & 1,798 & 400 & 400 & 2,598 \\
Inferolateral (\IL) & 450 & 100 & 100 & 650 \\
\bottomrule
\end{tabular}
\end{table}
\tracenote{Table 3.1 numbers: A§1.1.4 (label × split; folder × split) and A§1.1.7 (territory_4c counts); totals per territory computed as row sums}

\paragraph{The simulation parameters $\thetavec$.}
\label{sec:theta}
For each MI simulation the release ships the ischaemia parameters used to generate it. Four are read into the analysis: the angular position of the ischaemic region around the left ventricle, $\phiang \in (-\pi, \pi]$; its apico-basal position $z \in [0.10, 1.00]$; its size (75--175 in the simulator's units); and its transmural extent $\rho_{\varepsilon,\max} \in \{0.3, 1.0\}$. A fifth stored array, \texttt{transmural}, is a byte-identical duplicate of $\rho_{\varepsilon,\max}$ on all three splits and is not a separate parameter. \tracenote{A§1.1.6; data/theta_mi_*.npz; scripts/build_medalcare_isch_targets.py:71,308-311} Two properties of $\thetavec$ matter later. First, $z$ and size have identical ranges in every occlusion folder, so within the MI cohort the parameters that carry territory information are $(\phiang, \rho_{\varepsilon,\max})$ alone. Second, the four-class territory label used throughout, \territory, is defined from $\phiang$ by fixed wedges---$[0, 2]$\,rad $\to$ \AS, $(2, \pi] \to$ \AL, $[-2, 0) \to$ \INF, $[-\pi, -2) \to$ \IL---and not from the folder name. The two agree everywhere except in one folder, \texttt{LCX\_0.3\_post}, whose simulated angles fall in the anterolateral wedge; the build script relabels those rows and fails on any other disagreement. \tracenote{A§1.1.7; scripts/build_medalcare_isch_targets.py:143-168,350-387} Because \texttt{LCX\_0.3\_ant} and \texttt{LCX\_0.3\_post} occupy the same $\phiang$ wedge---the same distribution in all four parameters under two folder names---any readout of the \emph{folder} labels has a structural ceiling (macro-F$_1$ 0.864, with \IL\ recall pinned at 0.5). Results obtained before this was understood used the folder labels and were implicitly compared against 1.0; every four-class result in this thesis uses the $\phiang$-derived label, which has no such ceiling. \tracenote{A§1.1.7 :135-139; reports/2026-08-10_repo_audit_and_rerun_plan.md:140-153; dossier B §4.4}

The circular means of $\phiang$ within the four wedges---$+57.3^\circ$ (\AS), $+147.3^\circ$ (\AL), $-147.3^\circ$ (\IL), $-57.3^\circ$ (\INF)---serve as \emph{anchor angles} whenever a real-domain territory label has to be placed on the simulator's circle. They are a construction of the simulator's own labelling (they reproduce the wedge midpoints to $0.04^\circ$), not an independent measurement. \tracenote{A§1.1.8; outputs/analysis/circular_geometry/floor_audit.json→anchors_deg}

\subsection{PTB-XL: the real cohort}
\label{sec:data-ptbxl}

\ptbxl\ v1.0.3 \citep{wagner2020ptbxl,strodthoff2021deep} contains 21,799 ten-second 12-lead ECGs from 18,869 patients, distributed with an official ten-fold patient-stratified partition. The 500\,Hz high-resolution files are used, so both domains enter the encoder as $(12 \times 5000)$ arrays. Each record carries SCP diagnostic statements grouped into five superclasses (NORM, MI, STTC, HYP, CD); labels are built multi-hot over every listed statement without a likelihood threshold. \tracenote{A§1.2.1-1.2.2; scripts/datasets.py:551-709}

\paragraph{The shared three-class space.}
Fine-tuning that sees both domains needs a common label space. The one used from Experiment~7 onwards is $\{$NORM, MI, CD$\}$: on the synthetic side sinus $\to$ NORM, MI $\to$ MI, and the two bundle-branch blocks, interatrial block and AV block $\to$ CD, with left atrial enlargement and fibrotic atrial cardiomyopathy dropped; on the real side NORM and MI map to themselves and the CD superclass to CD, with STTC and HYP dropped for want of a clean correspondence. \tracenote{A§1.2.3; scripts/finetune_multilabel.py:50-84} Before the August audit the real-side filter mistakenly retained STTC in place of CD (defect~F, \S\ref{sec:preproc}); all pre-audit PTB-XL classification metrics were therefore computed on a different subset from the post-audit ones and the two are not comparable. \tracenote{A§7.7}

\paragraph{Infarct territory in the real cohort.}
Fourteen SCP statements belong to the MI superclass. They are mapped to the same four territories as the simulator by a fixed rule set: anterior-side codes (AMI, ASMI, INJAS) give \AS\ unless an anterolateral or lateral code (ALMI, INJAL, LMI, INJLA) is also present, in which case \AL; inferior-side codes (IMI, INJIN) give \INF\ unless an inferolateral or lateral code (ILMI, INJIL, IPLMI, LMI, INJLA) is present, in which case \IL; records with any posterior code (PMI, IPMI), with both anterior- and inferior-side codes, or with lateral codes only are excluded. \tracenote{A§1.2.4; scripts/build_ptbxl_mi_subclass.py:88-159} Table~\ref{tab:ptbxl} gives the resulting cohort. Two evaluation sets appear in the results: fold~10 alone ($n=438$), the only legitimate test set for encoders whose training saw folds 1--9; and all ten folds ($n=4{,}324$, 3,794 patients), usable only for the encoder that never took a PTB-XL gradient step (\S\ref{sec:encoder}). An older three-class rule (single-territory Anterior/Inferior/Lateral) gives 444 rows on fold~10; whenever a count is quoted the rule is named. \tracenote{A§1.2.5, A§7.10; data/ptbxl_mi_subclass_summary.json; data/ptbxl_mi_subclass_allfolds.csv}

\begin{table}[tb]
\centering\small
\caption{\ptbxl\ cohort. The four-class territory is derived from the MI-superclass SCP statements by the fixed rule in the text; the two-class collapse merges \AS+\AL\ (anterior) and \INF+\IL\ (inferior). Acuity is the \texttt{infarction\_stadium1} field where graded.}
\label{tab:ptbxl}
\begin{tabular}{lrr}
\toprule
 & fold 10 & all folds (1--10) \\
\midrule
records & 2,198 & 21,799 \\
MI superclass present & 550 & 5,469 \\
four-class territory assigned & 438 & 4,324 \\
\quad Anteroseptal / Anterolateral & 168 / 42 & 1,624 / 455 \\
\quad Inferior / Inferolateral & 196 / 32 & 1,914 / 331 \\
\quad two-class: anterior / inferior & 210 / 228 & 2,079 / 2,245 \\
unique patients (four-class rows) & --- & 3,794 \\
\midrule
\multicolumn{3}{l}{\emph{acuity among the 4,324 four-class rows}}\\
graded (\texttt{stadium1} $\neq$ unknown) & & 1,617 \\
\quad Stadium II--III or III & & 1,459 (90.2\,\%) \\
\quad Stadium I, I--II or II & & 158 \\
\bottomrule
\end{tabular}
\end{table}
\tracenote{Table 3.2: A§1.2.5 (both count tables) and A§1.2.6 (acuity); 158 = 94+3+61}

The acuity field is used only once, to test whether cross-domain transfer is better on acute records (Chapter~\ref{ch:results}); it is graded on 1,617 of the 4,324 rows and 90.2\,\% of the graded rows are Stadium II--III or III, i.e.\ the real MI cohort is dominated by chronic infarcts. \tracenote{A§1.2.6}

\subsection{Pre-processing, lead order and the August fixes}
\label{sec:preproc}

Both loaders return $(12 \times 5000)$ arrays at 500\,Hz in the standard order I, II, III, aVR, aVL, aVF, V1--V6, obtained by re-indexing the WFDB header's \texttt{sig\_name} rather than by position, and z-scored per lead. \tracenote{A§1.1.9, A§1.2.2, A§1.3} Neither convention held throughout the project. Until 10 August 2026 the synthetic loader declared its input order with aVL and aVF transposed and applied a single global scalar z-score, while the real loader re-indexed by name and normalised per lead; the synthetic domain therefore entered every pre-August experiment with its two inferior/high-lateral limb leads swapped and with a different amplitude convention. The transposition was found by reading the loader and confirmed on the raw signals through the limb-lead identities $\mathrm{aVL} = (\mathrm{I}-\mathrm{III})/2$ and $\mathrm{aVF} = (\mathrm{II}+\mathrm{III})/2$, which hold for the on-disk channel~4 as aVL and channel~5 as aVF; the identity check, not any downstream metric, is the evidence for the fix. \tracenote{A§1.3; reports/2026-08-10_lead_order_bug_diagnostic.md; dossier C Card A2} Two further defects were found in the same audit: the PTB-XL three-class filter kept STTC instead of CD, and predicted-probability columns were passed to the AUC routine in the classifier's internal class order, corrupting some thirty reported AUCs. All were fixed in the working tree, the pre-fix state was preserved at the git tag \texttt{pre-leadfix}, and every encoder used in Chapter~\ref{ch:results} (\texttt{exp8\_leadfix\_*}) was retrained after the fixes. Numbers from before the fixes appear only in Appendix~\ref{app:supp}, marked as such. \tracenote{reports/2026-08-10_repo_audit_and_rerun_plan.md §0, §2; dossier C Card A}

\paragraph{Territory as an angle.}
\label{sec:angle}
Because the simulator's territory is a continuous angle while the real cohort's is a four-way diagnosis, the two are compared on the circle by assigning each real record the anchor angle of its territory (\S\ref{sec:theta}). Every real-domain ``angle'' in this thesis is therefore a four-point quantisation of the simulator's labelling, not a measurement; this is what makes a constant-predictor floor for circular agreement unavoidable and is why \S\ref{sec:pipelines} defines one. \tracenote{A§7.11; analysis/geom_common.py:158-177}

\paragraph{Anatomy models and splits.}
The release's promise that an anatomical model never spans partitions holds for 11 of the 13 model directories; two (\texttt{run\_S64}, \texttt{run\_S67}) straddle a partition boundary. In-domain cross-validation in this thesis groups by simulation run, which is finer than the anatomical model, so this is recorded as a limitation (\S\ref{ch:conclusion}) rather than corrected. \tracenote{A§7.2}

\section{The foundation model and its adaptation}
\label{sec:encoder}
\todo{3.3 — ECGFounder Net1D, adapters (309,568 params) + head (3,075); freezing convention; modes single / joint dual-head (+51-name physics head $\theta_{\mathrm{phys}}$, distinct from $\thetavec$) / shared-head; MMD/ccMMD; bottleneck heads; Table 3.3 run families with pre/post-fix flag; why \texttt{exp8\_leadfix\_medalonly} is the headline (n=4324) and dual is sensitivity; latent export}

\section{The experimental programme, dated}
\label{sec:programme}
\todo{3.4 — Table 3.4 dated glossary (weekly baselines; Exp 1/4; Exp 5/6; Exp 7; Track 1; Track 2 INLP; Phase B/B2; Track 3; poster; audit + exp8; fair control / grid / sweep / probe map; circular geometry; F1–F3); Exp 2/3 never run; multi-seed never run; June–July gap}

\section{Hand-crafted ECG features}
\label{sec:features}
\todo{3.5 — global6 and spatial54 (NeuroKit2 dwt, lead-II fiducials, raw voltages); blocks; duplicates; coverage/MNAR; axis; the control arm and its instrument checks; Table 3.5}

\section{Analysis pipelines}
\label{sec:pipelines}
\todo{3.6 — P1 Phase-B2 classifier (LogReg, C by CV, scaler modes by code name: target = legacy, target\_pool\_measured = strict, source = defect; paired\_macro\_f1); P2 circular geometry (ridge to (cos,sin), GCV; \circR\ and the constant floor as supremum; permutation nulls below floor; floor-free metrics; block bootstrap; two mandatory nulls; axis baseline; scalers source = strict transport / target = re-centring); alignment diagnostics (C2ST linear/GBDT, MMD, kNN mixing, INLP, permutation sweep); Table 3.6 pipeline glossary; the self-corrections}

\section{The informativeness-fidelity audit}
\label{sec:audit-method}
\todo{3.7 — per-feature $\eta^2$ per domain; 500-draw block bootstrap; blind-spot / spurious criteria; block-aware permutation; |SMD|; 2-class collapse}

\section{Transport, prediction and repair protocols}
\label{sec:transport-method}
\todo{3.8 — readouts fit on MedalCare only; scalers named per pipeline (\qone{Q1 declaration}); efficiency; block readouts (F2); channel restriction via latent$\to$feature map (F3-A); reweighting + ESS (F3-B); diagonal CORAL boundary; paired tests; pre-statement}
